% Lean: EconLean/Macro/NK/Rotemberg.lean, EconLean/Finance/HeKrishnamurthy.lean
\chapter{Price Adjustment Costs and Intermediary Asset Pricing}
\label{chap:pricing_intermediaries}

Sources: Rotemberg (1982); Woodford, \emph{Interest and Prices} ch.~3; Gal\'i ch.~3; He--Krishnamurthy (2013).

\section{Rotemberg pricing}

\begin{definition}[Firm's quadratic loss]
\label{def:rot}
\lean{Econ.Rotemberg.loss}
\leanok
$\sum_{s<T}\beta^{s+1}[(q_{s+1}-p_{s+1}-mc_{s+1})^2+c(q_{s+1}-q_s)^2]$; the firm takes the aggregate path $p$ as given.
\end{definition}

\begin{theorem}[First-order condition]
\label{thm:rot_foc}
\lean{Econ.Rotemberg.hasDerivAt_update, Econ.Rotemberg.hasDerivAt_loss, Econ.Rotemberg.foc, Econ.Rotemberg.isLocalMin_of_foc}
\leanok
\uses{def:rot}
At an interior date $n=m+1<T$ the derivative of the loss in $q_n$ is
$2\beta^n[(q_n-p_n-mc_n)+c(q_n-q_{n-1})-\beta c(q_{n+1}-q_n)]$. A local minimum makes it vanish, and conversely (for
$c\ge0$) a zero makes $q_n$ a minimizer.
\end{theorem}
\begin{proof}
\leanok
Only the terms $s=m$ and $s=m+1$ depend on $q_n$. Sufficiency: the derivative is
$2\beta^n(1+c+\beta c)(x-q_n)$, so by the mean value theorem the loss rises in both directions.
\end{proof}

\begin{theorem}[Rotemberg NKPC and the Calvo bridge]
\label{thm:rot_nkpc}
\lean{Econ.Rotemberg.nkpc, Econ.Rotemberg.calvo_equiv, Econ.Rotemberg.calvo_bridge, Econ.Rotemberg.steady_inflation, Econ.Rotemberg.optimal_exists}
\leanok
\uses{thm:rot_foc}
In a symmetric equilibrium $\pi_n=\beta\pi_{n+1}+\frac1c mc_n$. With $c=\theta/((1-\theta)(1-\beta\theta))$ this is
exactly the Calvo NKPC. Steady inflation $p_n=n\bar\pi$ with $mc=c(1-\beta)\bar\pi$ is an equilibrium.
\end{theorem}
\begin{proof}
\leanok
Set $q=p$ in the first-order condition and divide by $c$; the slope identity is algebra; the witness satisfies the
first-order condition, which is sufficient.
\end{proof}

\section{Intermediary asset pricing}

\begin{definition}[Equity-constrained intermediary]
\label{def:hk}
\lean{Econ.HeKrishnamurthy.Env, Econ.HeKrishnamurthy.Env.capital, Econ.HeKrishnamurthy.Env.obj, Econ.HeKrishnamurthy.Env.premium, Econ.HeKrishnamurthy.Env.IsEq}
\leanok
Capital $K=w+\min(W,mw)$; holding $x$ yields $x\mu-\gamma\sigma^2x^2/(2K)$; equilibrium: holding the unit supply is
optimal.
\end{definition}

\begin{theorem}[Risk premium]
\label{thm:hk_premium}
\lean{Econ.HeKrishnamurthy.Env.capital_pos, Econ.HeKrishnamurthy.Env.isMax_iff, Econ.HeKrishnamurthy.Env.isEq_iff, Econ.HeKrishnamurthy.Env.premium_constrained, Econ.HeKrishnamurthy.Env.premium_unconstrained, Econ.HeKrishnamurthy.Env.sharpe_eq, Econ.HeKrishnamurthy.Env.exists_eq, Econ.HeKrishnamurthy.Env.both_regions}
\leanok
\uses{def:hk}
The optimal holding is $K\mu/(\gamma\sigma^2)$, so $\mu=\gamma\sigma^2/K$: $\gamma\sigma^2/((1+m)w)$ when
$mw<W$, and $\gamma\sigma^2/(w+W)$ otherwise.
\end{theorem}
\begin{proof}
\leanok
Complete the square: the objective is its maximum minus $\gamma\sigma^2(x-x^*)^2/(2K)$.
\end{proof}

\begin{theorem}[Intermediary wealth as a state variable]
\label{thm:hk_state}
\lean{Econ.HeKrishnamurthy.Env.capital_strictMono, Econ.HeKrishnamurthy.Env.premium_strictAnti, Econ.HeKrishnamurthy.Env.constrained_elasticity, Econ.HeKrishnamurthy.Env.unconstrained_redistribution, Econ.HeKrishnamurthy.Env.constrained_redistribution_matters}
\leanok
\uses{thm:hk_premium}
The premium falls strictly with intermediary wealth. When constrained it has unit elasticity, and shifting wealth to
intermediaries at fixed total wealth lowers it; when unconstrained only total wealth matters.
\end{theorem}
\begin{proof}
\leanok
Capital is strictly increasing in $w$; the rest follows from the two closed forms.
\end{proof}
